% IEEE conference paper. Starter file written for OpenPaperwork; it uses the IEEEtran class
% that ships with the TeX distribution and copies nothing from IEEE's own template files.
\documentclass[conference]{IEEEtran}
\usepackage{cite}
\usepackage{amsmath}
\usepackage{graphicx}
\usepackage{booktabs}
\usepackage[hidelinks]{hyperref}

\begin{document}

\title{Cold-Start Latency in Managed Vector Search:\\A Two-Week Measurement}

% For blind review keep this block exactly as it is. For the final version, replace it with
% \IEEEauthorblockN{Name} and \IEEEauthorblockA{Department, Organisation \\ City, Country \\ email}
\author{\IEEEauthorblockN{Anonymous Author(s)}
\IEEEauthorblockA{Submitted for blind review}}

\maketitle

\begin{abstract}
Managed vector indexes are usually described by their steady-state recall and throughput.
We instead measure the first query after an idle period, the cold start, on one managed
service over fourteen days. Across 2{,}480 cold starts the median latency was 780\,ms,
seven times the warm median, and the slowest decile exceeded four seconds. Cold-start cost
scaled with index size but not with the number of replicas, which suggests the delay sits
in loading index segments rather than in request routing. We report the method, the
distribution, and the two configuration changes that removed most of the tail.
\end{abstract}

\begin{IEEEkeywords}
vector search, cold start, latency measurement, managed services
\end{IEEEkeywords}

\section{Introduction}
Approximate nearest-neighbour indexes are memory resident by design. When a managed
service scales a replica down, the next query must rebuild or reload the index before it
can answer. Operators see this as an occasional slow query; users see it as a first
interaction that feels broken. Published evaluations report warm percentiles, because
warm measurement is easy to automate \cite{example-survey}.

This paper asks a narrow question: how large is the cold-start penalty in practice, and
what does it scale with? We measured one managed vector service for fourteen days,
issuing the first query of every idle window and recording the service's own latency
figures alongside client-side timing.

\section{Method}
We issued one query every fifteen minutes to a three-replica index of 4.2 million
768-dimensional vectors, and treated a query as cold when the preceding idle gap exceeded
ten minutes. Warm queries ran in a separate stream against the same index so that both
distributions come from the same days and the same data. We recorded client time to first
byte, service-reported latency, and the replica that answered.

\begin{table}[b]
\caption{Latency by state, milliseconds, over fourteen days}
\label{tab:latency}
\centering
\begin{tabular}{lrrr}
\toprule
State & Median & p90 & Max \\
\midrule
Warm        & 110 & 190  & 640 \\
Cold        & 780 & 4{,}100 & 9{,}800 \\
\bottomrule
\end{tabular}
\end{table}

\section{Results}
Table~\ref{tab:latency} gives the two distributions. The cold median is 7.1 times the warm
median, and the gap widens in the tail: the cold p90 is more than twenty times the warm
p90. Cold starts were 3.1 per cent of queries but 41 per cent of the total time spent
waiting.

Cold-start latency correlated with index size ($r = 0.71$) and not with replica count
($r = 0.08$), which points at segment loading rather than scheduling. Raising the minimum
replica count from one to two, and enabling the service's keep-warm option, removed 86 per
cent of cold starts above one second at a cost of 19 per cent more instance hours.

\section{Discussion}
The measurement is limited to one service, one index shape and one region, so the numbers
are not a ranking. What generalises is the shape: a heavy tail that steady-state
benchmarks do not show and that a keep-warm setting can remove. We recommend that managed
services report a cold-start percentile beside their warm figures, since the first query
of a session is the one a person remembers.

\begin{thebibliography}{2}
\bibitem{example-survey}
A.~Example and B.~Sample, ``A survey of approximate nearest-neighbour benchmarks,''
\emph{Journal of Invented Systems}, vol.~12, no.~3, pp.~210--219, 2025.

\bibitem{example-notes}
C.~Placeholder, ``Notes on index reload behaviour in shared clusters,'' in \emph{Proc.
Invented Conference on Systems Measurement}, 2026, pp.~44--51.
\end{thebibliography}

\end{document}
